\documentclass[11pt]{article}

\usepackage[T1]{fontenc}
\usepackage[utf8]{inputenc}
\usepackage{lmodern}
\usepackage[a4paper,margin=28mm]{geometry}
\usepackage{microtype}
\usepackage{amsmath,amssymb,amsthm,mathtools}
\usepackage{enumitem}
\usepackage{booktabs,tabularx}
\usepackage[colorlinks=true,linkcolor=blue,citecolor=blue,urlcolor=blue]{hyperref}

\newtheorem{theorem}{Theorem}[section]
\newtheorem{proposition}[theorem]{Proposition}
\newtheorem{corollary}[theorem]{Corollary}
\theoremstyle{definition}
\newtheorem{definition}[theorem]{Definition}
\newtheorem{example}[theorem]{Example}
\theoremstyle{remark}
\newtheorem{remark}[theorem]{Remark}

\newcommand{\vol}{\operatorname{vol}}
\newcommand{\card}{\operatorname{card}}

\title{Causal Sets as a Conditional Coarse-Graining of Modal Triplet Theory:\\
\large Kinematic Construction, Sampling Choices, and Scale-Typing Boundary}
\author{Peter Nero}
\date{July 2026\\Version 2}

\begin{document}
\maketitle

\begin{abstract}
A causal set can be obtained from an effective Modal Triplet Theory
spacetime once a sampling rule is supplied, but the sampling rule is
additional data. For a globally hyperbolic four-dimensional spacetime, any
locally finite set of sampled events inherits a partial order from the
spacetime causal relation. A Poisson process with intensity proportional to
the spacetime volume measure is almost surely locally finite on every
relatively compact region; in Minkowski spacetime its law is Lorentz
invariant because the intensity measure is invariant. These are exact
kinematic statements. They do not identify the causal set as fundamental,
derive the sprinkling intensity, or convert an internal q79 spectral gap
into a four-dimensional event density. Such a conversion requires a
base-resolution theorem with metric normalization and an explicit
internal-to-spacetime map. Causal sets therefore remain a valid conditional
coarse-graining or observational encoding of an MTT spacetime, not a
currently selected consequence of the internal finite geometry.
\end{abstract}

\section*{Version 2 Revision Note}
\begin{description}[leftmargin=3.2cm,style=nextline]
\item[Supersedes] Version 1.0, \emph{Causal Sets as an Effective Limit of
Modal Triplet Theory}.
\item[Reason] The original paper identified an internal modal gap with a
four-dimensional sprinkling density and described Poisson sampling as
derived rather than chosen.
\item[Resolution] The spacetime, point process, intensity measure, causal
order, and local-finiteness assumptions are now typed separately. The exact
kinematic construction is proved without a cross-sector scale claim.
\item[Retained result] A locally finite sample of a causal spacetime inherits
a causal-set order, and covariant Poisson sprinkling gives the standard
Lorentz-invariant kinematic model in Minkowski spacetime.
\item[Remaining boundary] MTT must derive a physical event/sampling law and
its four-dimensional intensity from the selected branch before the causal
set can be called an MTT prediction.
\end{description}

\tableofcontents

\section{Causal sets and effective spacetime}

A causal set is a pair \((C,\preceq)\) such that:
\begin{enumerate}
\item \(\preceq\) is reflexive;
\item \(\preceq\) is antisymmetric;
\item \(\preceq\) is transitive;
\item every order interval
\[
 I(x,y)=\{z\in C:x\preceq z\preceq y\}
\]
is finite.
\end{enumerate}
The order represents causal precedence, while local finiteness supplies
discreteness \cite{Bombelli,Sorkin}.

MTT currently reaches a selected globally hyperbolic four-dimensional
spacetime at a conditional gravity/QFT tier. The question in this paper is
not whether spacetime is fundamentally discrete. It is whether a causal-set
description can be extracted from that effective spacetime.

\section{The exact kinematic construction}

\begin{theorem}[Inherited causal-set order]
\label{thm:inherit}
Let \((Y_4,g)\) be a causal spacetime, and let \(C\subset Y_4\) be a set such
that \(C\cap K\) is finite for every compact \(K\subset Y_4\). Define
\[
 x\preceq y
 \quad\Longleftrightarrow\quad
 y\in J^+(x).
\]
If every causal diamond \(J^+(x)\cap J^-(y)\) is compact, then
\((C,\preceq)\) is a causal set.
\end{theorem}

\begin{proof}
The causal relation is reflexive and transitive. Causality excludes a
nontrivial closed causal curve, giving antisymmetry. For \(x\preceq y\),
\[
 I(x,y)
 =
 C\cap J^+(x)\cap J^-(y).
\]
The diamond is compact, so the assumed local finiteness of \(C\) makes this
intersection finite.
\end{proof}

Global hyperbolicity supplies the compact-diamond hypothesis. The theorem is
purely kinematic: it uses a continuum causal spacetime and a locally finite
sample. It does not determine how the sample is generated.

\section{Poisson sprinkling}

Let \(\mu_g\) be the spacetime volume measure and let \(\rho>0\). A Poisson
point process with intensity measure
\[
 \Lambda(B)=\rho\,\mu_g(B)
\]
has
\[
 \Pr[N(B)=n]
 =
 e^{-\Lambda(B)}\frac{\Lambda(B)^n}{n!}
\]
for measurable regions of finite volume, with independent counts on
disjoint regions.

\begin{theorem}[Poisson local finiteness]
\label{thm:poisson}
On a sigma-finite spacetime measure space, a Poisson process with locally
finite intensity measure has finitely many points in every relatively
compact finite-volume region almost surely. Combined with
\autoref{thm:inherit} on a globally hyperbolic spacetime, it produces a
causal set almost surely.
\end{theorem}

\begin{proof}
For such a region \(B\), \(N(B)\) is a Poisson random variable with finite
mean \(\Lambda(B)\), and therefore takes a finite integer value almost
surely. A countable exhaustion by relatively compact regions gives local
finiteness simultaneously.
\end{proof}

\subsection{Lorentz invariance}

In Minkowski spacetime, the volume measure is invariant under proper
orthochronous Poincare transformations. Therefore the law of a homogeneous
Poisson process with intensity \(\rho\,d^4x\) is invariant. Individual
realizations are not symmetric configurations; the probability law is.

On a generic curved spacetime the analogous statement is covariance under
isometries of \((Y_4,g)\), not global Lorentz invariance. A deterministic
lattice or foliation-based sampling can select a preferred frame.

\section{Why the density is not derived from an internal gap}

Suppose \(\lambda_{\mathrm{int}}\) is an eigenvalue of an operator on the
internal compactification \(X_6\). A four-dimensional sprinkling intensity
has units of inverse four-volume. The relation
\[
 \rho\stackrel{?}{=}f(\lambda_{\mathrm{int}})
\]
is not defined until the internal metric normalization, four-dimensional
metric normalization, reduction moduli, and a cross-sector response map are
fixed.

\begin{proposition}[Scale-typing obstruction]
\label{prop:scale}
An internal spectral gap alone cannot canonically determine a
four-dimensional sprinkling intensity.
\end{proposition}

\begin{proof}
Rescale the internal and external metrics independently. The internal
eigenvalue and the external four-volume transform with independent powers.
No relation between them is invariant unless additional reduction data
constrain the rescalings and define a unit-preserving map.
\end{proof}

The same obstruction appears in the finite-filter and gravitational-collapse
papers. It does not say that a relation is impossible. It identifies the
missing theorem.

\section{Alternative coarse-grainings}

Poisson sprinkling is not the only way to produce a locally finite sample.
Examples include:
\begin{itemize}
\item detector records in a bounded observational protocol;
\item a covariantly defined point process with non-Poisson correlations;
\item a finite event set produced by a hybrid selection dynamics;
\item an adaptive cover or net used only for numerical approximation.
\end{itemize}

Each choice answers a different question. Detector records describe an
operational history. A Poisson process is a statistically Lorentz-invariant
encoding. A numerical net is a discretization tool. None should be called
the fundamental event ontology without an additional physical argument.

The companion event-selection paper gives sufficient dwell-time and hazard
conditions for dynamically generated events to be locally finite. That
dynamic problem is deliberately not duplicated here.

\section{What is and is not reconstructed}

The order plus number principle of causal-set theory suggests that causal
order and counting density can encode much of Lorentzian geometry under
appropriate manifoldlike conditions. The present construction starts from
the continuum geometry, so it does not prove the reverse reconstruction or
manifoldlikeness of an arbitrary causal set.

\begin{center}
\begin{tabularx}{\textwidth}{@{}
>{\raggedright\arraybackslash}p{0.31\textwidth}
>{\raggedright\arraybackslash}p{0.18\textwidth}
>{\raggedright\arraybackslash}X@{}}
\toprule
Statement & Status & Boundary\\
\midrule
Locally finite sample inherits causal order & Exact &
\autoref{thm:inherit}.\\
Poisson sprinkling is locally finite & Exact &
\autoref{thm:poisson}.\\
Minkowski Poisson law is Lorentz invariant & Exact &
Invariance is statistical, not realization-wise.\\
Causal set reconstructs the input continuum & Not shown &
Requires manifoldlikeness and reconstruction results.\\
MTT selects Poisson statistics & Open &
No selected event or sampling law.\\
MTT selects the density from q79 gap & Open/ill-typed without bridge &
Needs base-resolution and unit map.\\
Causal set is fundamental in MTT & Not claimed &
Current role is effective encoding.\\
\bottomrule
\end{tabularx}
\end{center}

\section{Discussion}

The corrected relation between MTT and causal sets is straightforward:
\[
\text{selected effective }(Y_4,g)
\;+\;
\text{declared locally finite sampling law}
\longrightarrow
\text{causal set}.
\]
That result is useful. It permits causal-set observables and reconstruction
tools to be applied to an MTT spacetime, and it provides a concrete way to
compare continuum and discrete descriptions.

The missing density theorem is also a productive target. A valid derivation
would map selected q79 geometry and the four-dimensional metric to an event
intensity or detector-resolution functional with units and covariance
preserved. Until that map exists, varying \(\rho\) is an encoding parameter,
not a prediction.

\section{Conclusion}

Causal sets are compatible with MTT as a conditional coarse-graining of its
effective spacetime sector. The inherited order and Poisson local-finiteness
results are exact. The physical sampling law and its density are not.

This distinction preserves both programs. Causal-set mathematics can be used
without claiming that MTT has derived fundamental discreteness, and MTT can
seek a dynamical event law without building Poisson statistics into its
premises.

\begin{thebibliography}{99}

\bibitem{Bombelli}
L.~Bombelli, J.~Lee, D.~Meyer, and R.~D.~Sorkin,
``Space-time as a causal set,''
\emph{Physical Review Letters} \textbf{59} (1987), 521--524.

\bibitem{Sorkin}
R.~D.~Sorkin,
``Causal sets: Discrete gravity,''
in A.~Gomberoff and D.~Marolf, eds.,
\emph{Lectures on Quantum Gravity},
Springer, 2005, pp.~305--327.

\bibitem{BrightwellGregory}
G.~Brightwell and R.~Gregory,
``Structure of random discrete spacetime,''
\emph{Physical Review Letters} \textbf{66} (1991), 260--263.

\bibitem{DaleyVereJones}
D.~J.~Daley and D.~Vere-Jones,
\emph{An Introduction to the Theory of Point Processes, Volume I},
Springer, 2003.

\bibitem{MTTFoundation}
P.~Nero,
\emph{Modal Triplet Theory: Foundations},
current revised MTT paper corpus, 2026.

\bibitem{MTTGR}
P.~Nero,
\emph{Controlled Coherent Reduction to Four-Dimensional Einstein Gravity},
current revised MTT paper corpus, 2026.

\end{thebibliography}

\end{document}
